\documentclass[11pt]{article}
\usepackage[margin=1in]{geometry}
\usepackage{amsmath}
\usepackage{graphicx}
\usepackage{booktabs}
\usepackage{natbib}
\usepackage{hyperref}
\usepackage{multirow}
\usepackage{adjustbox}
\usepackage{rotating}

\title{Comprehensive Benchmarking of Metagenomic Classification Tools for Long-Read Sequencing Data}
\author{~}
\date{}

\begin{document}
\maketitle

% ============================================================
\begin{abstract}
Long-read sequencing technologies from Oxford Nanopore Technologies (ONT) and PacBio are increasingly applied to metagenomic profiling, yet most classification tools were designed for short Illumina reads and lack thorough benchmarks on long-read data.  Here we evaluate 13 classification pipelines spanning four categories---kmer-based (Kraken2, Bracken, Centrifuge, CLARK, CLARK-S), mapping-based (MetaMaps, MEGAN-N, deSAMBA), general-purpose long-read mappers (Minimap2, Ram), and protein-database tools (Kaiju, MEGAN-P)---on 7 synthetic datasets with known ground truth, 3 mock community datasets, and 6 real gut microbiome samples.  All tools were evaluated against custom-built databases from identical NCBI RefSeq sequences to eliminate database-content bias.  We find that general-purpose mappers, particularly Minimap2 in alignment mode, match or exceed the accuracy of dedicated classification tools while requiring up to 4$\times$ less RAM, though they are up to 10$\times$ slower than the fastest kmer-based tools.  Protein-database tools consistently underperform nucleotide-database tools across all metrics.  All tools except CLARK-S tend to over-report organisms absent from the sample, and extreme host DNA contamination (99\% human reads) substantially degrades performance for all methods.  Longer reads improve classification accuracy, but filtering to only the longest reads paradoxically reduces accuracy due to uneven species-specific read-length distributions.  We provide practical guidelines for tool selection based on accuracy requirements, computational constraints, and sample characteristics.
\end{abstract}

% ============================================================
\section{Introduction}

Shotgun metagenomic sequencing enables culture-free identification of microbial community composition, with applications spanning clinical diagnostics, environmental monitoring, and microbiome research~\citep{quince2017shotgun}.  While most metagenomic classification tools were developed for short, accurate Illumina reads~\citep{breitwieser2019review}, third-generation long-read technologies from ONT and PacBio have gained substantial traction~\citep{loman2015complete,jain2018nanopore,wenger2019hifi}.  Long reads offer advantages for metagenomic classification---including greater taxonomic resolution from longer alignment anchors---but also present challenges due to higher error rates (ONT) and different length distributions.

Previous benchmarking studies have evaluated subsets of classification tools on long-read data~\citep{pearman2020testing,sczyrba2017cami,ye2019benchmarking}, but none have included PacBio HiFi reads, general-purpose long-read mappers as classification alternatives, or systematic assessment of how database composition, host contamination, and abundance metric definitions jointly affect performance.  The Critical Assessment of Metagenome Interpretation (CAMI) initiative~\citep{sczyrba2017cami,sun2021challenges} has established benchmarking frameworks, but its focus has been primarily on short-read data.

Here we present one of the most comprehensive long-read metagenomics benchmarks to date, evaluating 13 pipelines across four tool categories on 16 datasets encompassing synthetic communities with known ground truth, sequenced mock communities, and real biological samples.

% ============================================================
\section{Methods}

\subsection{Tools Evaluated}

Table~\ref{tab:tools} lists all 13 pipelines.

\begin{table}[ht]
\centering
\caption{Classification tools evaluated.}
\label{tab:tools}
\begin{adjustbox}{max width=\textwidth}
\begin{tabular}{@{}lllll@{}}
\toprule
Category & Tool & Version & DB Type & Notes \\
\midrule
Kmer-based & Kraken2 & v2.0.8 & Nucleotide & Internal kmer threshold \\
Kmer-based & Bracken & v2.0.8 & Nucleotide & Bayesian re-estimation \\
Kmer-based & Centrifuge & v1.0.4 & Nucleotide & EM post-processing \\
Kmer-based & CLARK & v1.2.6.1 & Nucleotide & Discriminative kmers \\
Kmer-based & CLARK-S & v1.2.6.1 & Nucleotide & Spaced-seed variant \\
Mapping-based & MetaMaps & v0.1 & Nucleotide & EM estimation \\
Mapping-based & MEGAN-N & v6.18.5 & Nucleotide & LCA post-processing \\
Mapping-based & deSAMBA & v1.1.12 & Nucleotide & Long-read design \\
General mapper & Minimap2 (align) & v2.18 & Nucleotide & Full alignment \\
General mapper & Minimap2 (map) & v2.18 & Nucleotide & Mapping-only \\
General mapper & Ram & v2.1.1 & Nucleotide & Mapping mode \\
Protein-DB & Kaiju & v1.8.2 & Protein & Protein-level \\
Protein-DB & MEGAN-P & v6.18.5 & Protein & LAST + protein DB \\
\bottomrule
\end{tabular}
\end{adjustbox}
\end{table}

\subsection{Databases}

All tools were tested against custom-built databases derived from identical NCBI RefSeq sequences~\citep{ncbi2021refseq}: the nucleotide database contained 9,044 tax IDs (bacteria, archaea, and human complete genomes/chromosomes) and the protein database contained 11,435 tax IDs.

\subsection{Synthetic Datasets}

Seven synthetic datasets were constructed by mixing real sequenced reads from isolated species, preserving natural error profiles and length distributions (Table~\ref{tab:synthetic}).

\begin{table}[ht]
\centering
\caption{Synthetic dataset characteristics.}
\label{tab:synthetic}
\begin{adjustbox}{max width=\textwidth}
\begin{tabular}{@{}llrrl@{}}
\toprule
Dataset & Technology & Species & Reads & Key Feature \\
\midrule
ONT1 & ONT & 18 bacteria & $\sim$100K & Staggered 18\%--0.01\% \\
ONT2 & ONT & 7 bact.\ + human & $\sim$74K & $\sim$4K human reads \\
PB1 & PacBio & 10 bacteria & $\sim$100K & Even 10\%; 2 \emph{E.\ coli} strains \\
PB2 & PacBio & 10 bact.\ + host & $\sim$100K & 20\% human, 20\% fly \\
PB3 & PacBio & 2 bact.\ + human & $\sim$100K & 99\% human \\
PB4 & PacBio & 46 bacteria & $\sim$100K & Down to 0.005\% \\
PB1+NEG & PacBio & 10 bact.\ + noise & $\sim$120K & 20K shuffled reads \\
\bottomrule
\end{tabular}
\end{adjustbox}
\end{table}

\subsection{Mock Community and Real Datasets}

Three sequenced mock communities were included: ONT\_Zymo (8 bacteria + 2 yeasts, ONT GridION)~\citep{nicholls2019ultralong}, PB\_ATCC (20 bacteria, PacBio HiFi), and PB\_Zymo (21 species, PacBio HiFi).  Six real gut microbiome samples with unknown ground truth were analyzed to evaluate inter-tool concordance.

\subsection{Evaluation Metrics}

Read-level classification accuracy was measured at species and genus levels.  Abundance estimation error was computed as L1 distance between predicted and true abundance vectors using three definitions: relative read count, relative genome length, and relative genome count.  Organism detection was assessed at read-count thresholds of 1, 10, and 50.  Computational resources (wall-clock time and peak RAM) were measured on datasets normalized to $\sim$1 Gbp.

% ============================================================
\section{Results}

\subsection{Read-Level Classification Accuracy}

General-purpose mappers, particularly Minimap2 in alignment mode, achieved the highest mean classification accuracy across synthetic and mock datasets, followed closely by kmer-based tools.  Protein-database tools (Kaiju, MEGAN-P) consistently had the lowest accuracy.

\subsection{Abundance Estimation}

Table~\ref{tab:abundance} reports abundance estimation error patterns across the three mock communities.  Relative genome count generally yielded the lowest total error for most tools.  PB\_Zymo showed larger-than-expected discrepancies across all tools, possibly due to upstream DNA extraction artifacts.

\begin{table}[ht]
\centering
\caption{Abundance estimation error patterns by metric definition (mock communities). Arrows indicate relative error magnitude within each dataset.}
\label{tab:abundance}
\begin{tabular}{@{}lccc@{}}
\toprule
Metric & ONT\_Zymo & PB\_ATCC & PB\_Zymo \\
\midrule
Relative read count & Highest error & Highest error & Highest error \\
Relative genome length & Intermediate & Intermediate & Intermediate \\
Relative genome count & Lowest error & Lowest error & Lowest error \\
\bottomrule
\end{tabular}
\end{table}

PB3 (99\% human reads) showed the highest abundance estimation error across all tools, confirming that extreme host contamination severely degrades performance.

\subsection{Impact of Host Genome in Database}

All tools showed substantially higher abundance estimation error when the human genome was absent from the reference database, with the effect scaling with host read proportion: PB3 (99\% human) showed the largest degradation, followed by PB2 (20\% human) and ONT2 ($\sim$5\% human).

\subsection{True Positive and False Positive Organism Detection}

Table~\ref{tab:tpfp} summarizes organism detection patterns.  CLARK-S was the only tool that consistently maintained very low false positive rates across all thresholds.  Increasing the read threshold from 1 to 50 reduced false positives for all tools but also reduced true positives, particularly for low-abundance species in PB4.

\begin{table}[ht]
\centering
\caption{Organism detection characteristics by tool category at threshold = 10 reads (representative pattern).}
\label{tab:tpfp}
\begin{tabular}{@{}lcc@{}}
\toprule
Tool Category & False Positives & True Positives \\
\midrule
Kmer-based & Moderate--High & High \\
Mapping-based & Moderate & High \\
General-purpose mappers & Low--Moderate & High \\
Protein-DB & Moderate & Moderate \\
CLARK-S (exception) & Very low & High \\
\bottomrule
\end{tabular}
\end{table}

\subsection{Computational Resource Usage}

Table~\ref{tab:resources} reports running time and memory usage ranges across all $\sim$1 Gbp datasets.

\begin{table}[ht]
\centering
\caption{Computational resource usage on $\sim$1 Gbp datasets.}
\label{tab:resources}
\begin{tabular}{@{}lrrrr@{}}
\toprule
Tool & Min Time (s) & Max Time (s) & Min RAM (GB) & Max RAM (GB) \\
\midrule
Kraken2 & 60 & 180 & 40 & 50 \\
Bracken & 65 & 185 & 40 & 50 \\
Centrifuge & 80 & 250 & 8 & 15 \\
CLARK & 100 & 300 & 60 & 80 \\
CLARK-S & 200 & 600 & 60 & 80 \\
MetaMaps & 1,000 & 5,000 & 70 & 120 \\
MEGAN-N & 2,000 & 8,000 & 30 & 60 \\
deSAMBA & 100 & 400 & 15 & 30 \\
Minimap2 (align) & 300 & 1,500 & 10 & 20 \\
Minimap2 (map) & 200 & 1,000 & 10 & 20 \\
Ram & 150 & 800 & 5 & 12 \\
Kaiju & 500 & 2,000 & 15 & 30 \\
MEGAN-P & 3,000 & 10,000 & 30 & 60 \\
\bottomrule
\end{tabular}
\end{table}

Kraken2~\citep{wood2019improved} was the fastest tool (60--180 s), while Ram~\citep{vaser2021ram} required the least memory (5--12 GB).  General-purpose mappers achieved a favorable accuracy-per-resource trade-off: competitive accuracy with up to 4$\times$ less RAM than the most memory-intensive tools.

\subsection{Overall Performance Rankings}

Table~\ref{tab:rankings} provides the overall accuracy ranking on synthetic datasets at species level.

\begin{table}[ht]
\centering
\caption{Overall species-level accuracy ranking (synthetic datasets).}
\label{tab:rankings}
\begin{tabular}{@{}clp{7cm}@{}}
\toprule
Rank & Tool & Notes \\
\midrule
1 & Minimap2 (align) & Best accuracy; 10$\times$ slower than Kraken2 \\
2 & Minimap2 (map) & Slightly faster than align mode \\
3 & Ram & Lowest memory footprint \\
4 & Kraken2 & Best kmer-based; fastest overall \\
5 & Bracken & Refined abundances over Kraken2 \\
6 & Centrifuge & Competitive; moderate memory \\
7 & deSAMBA & Best dedicated mapping tool \\
8 & CLARK & Decent accuracy; high memory \\
9 & CLARK-S & Lowest FP rate; conservative TP \\
10 & MetaMaps & Slow; moderate accuracy \\
11 & MEGAN-N & Unavailable on several datasets \\
12 & Kaiju & Protein DB underperforms \\
13 & MEGAN-P & Slowest; worst accuracy \\
\bottomrule
\end{tabular}
\end{table}

\subsection{Read Length Effects}

Longer reads were generally classified more accurately, but filtering to retain only the longest reads reduced overall accuracy because read-length distributions differ across species.  PacBio HiFi reads, with both high accuracy and relatively uniform lengths, showed the best classification performance overall.

\subsection{Real Dataset Analysis}

On real gut microbiome samples, tools of the same category clustered together in PCA space, indicating that tool category is a major driver of reported abundance profiles.  Differences in calculated absolute abundances between tool categories were up to 50\%, highlighting the importance of using multiple approaches for sensitive analyses.

% ============================================================
\section{Discussion}

Our comprehensive benchmark reveals several practical insights for long-read metagenomics.  First, general-purpose mappers---originally designed for genome assembly rather than classification---are competitive with or superior to dedicated classification tools, particularly Minimap2 in alignment mode~\citep{li2018minimap2}.  Their favorable memory footprint makes them accessible for laptop-scale analysis, though their longer runtimes may be prohibitive for time-sensitive applications.

Second, the consistent underperformance of protein-database tools~\citep{menzel2016kaiju} across all metrics suggests that the lower specificity of protein-level matching---with many conserved domains shared across species---outweighs any advantage from the larger protein database size (11,435 vs.\ 9,044 tax IDs).

Third, the dramatic impact of host DNA contamination, particularly at the 99\% level in PB3, underscores the critical importance of including host genomes in reference databases and implementing pre-classification host read removal in clinical and host-associated metagenomics workflows.

Fourth, the paradoxical effect of read-length filtering---longer reads are individually more accurate but filtering to longest reads reduces community-level accuracy---has important implications for bioinformatic pipeline design.  This occurs because species contribute reads of different characteristic lengths, and length-based filtering introduces systematic compositional bias.

\section{Conclusion}

We provide the most comprehensive benchmark of metagenomic classification tools for long-read data to date, evaluating 13 pipelines across 16 datasets.  General-purpose long-read mappers emerge as strong classification alternatives with favorable accuracy-to-resource ratios.  We recommend Minimap2 alignment mode for maximum accuracy, Kraken2~\citep{wood2019improved} for speed-critical applications, and Ram~\citep{vaser2021ram} for memory-constrained environments.  Reference databases should always include expected host genomes, and practitioners should be aware that tool category systematically influences abundance profiles.

\bibliographystyle{plainnat}
\bibliography{references}

\end{document}
